\section{Implementation details}
\label{app:implementation}

\budgetactor{}, AWM, and ASI are all run using the unified codebase and prompts released by \citet{wang-etal-2025-asi} to ensure that differences in results cannot be attributed to inconsistencies in agent scaffolding, prompt formatting, or environment interaction.
The codebase uses BrowserGym~\citep{chezelles-etal-2025-browsergym} for interacting with the environment.
AWM and ASI come from the same research group~\citep{wang-etal-2024-agent, wang-etal-2025-asi}, and AWM is included as a baseline in the more recent ASI codebase.
ReasoningBank is run using the codebase released by \citet{ouyang-etal-2025-reasoningbank} due to its distinct prompt structure and agent logic.

We set the actor LLM temperature of \budgetactor{} to $0$, the same as the default value in AWM and ASI.
ReasoningBank sets the temperature to $0.7$ and we kept the default intact.
The auxiliary modules of AWM, ASI, and ReasoningBank use their respective defaults, which we do not modify.
ReasoningBank's codebase also ships a vendored WebArena with a custom evaluation pipeline that replaces the standard BrowserGym harness.
All methods in this paper are evaluated using the unmodified BrowserGym evaluation harness to ensure identical assessment criteria across configurations.

We host Qwen~3.6-27B on four NVIDIA H100 GPUs using vLLM~\citep{kwon-etal-2023-vllm}. 
We used AI assistants during implementation and manuscript preparation for code-related tasks, including generating plotting scripts and LaTeX snippets.
We also used them to assist with writing, and for polishing and shortening paragraphs.


\section{Accessibility-tree pruning}
\label{app:pruning}

The accessibility-tree simplification used in \budgetactor{} is a deterministic rule-based function with no LLM call, model fine-tuning, or learned component, and therefore no additional token cost.
The function operates on the BrowserGym-style accessibility-tree, in which each node has a role (\textit{e.g.}, \texttt{button}, \texttt{StaticText}, \texttt{heading}) and an accessible name, and is rendered as one indented line per node.

\paragraph{Pruning rule.}
A \texttt{StaticText} child node is removed when its text content is already contained verbatim in the accessible name of its immediate parent.
This handles a common pattern in web accessibility-trees, where a parent node (such as a labelled button or a list item) carries the same text as its \texttt{StaticText} child.
Removing the duplicate cuts context length without dropping information.
\texttt{StaticText} nodes whose stripped content is empty are also removed.
Non-\texttt{StaticText} nodes, node IDs (BIDs), roles, and tree structure are not modified.
% 
The rule has one exception: a \texttt{StaticText} child is kept whenever its parent's role is one of \{\texttt{article}, \texttt{paragraph}, \texttt{heading}, \texttt{strong}, \texttt{emphasis}, \texttt{mark}, \texttt{sectionheader}\}, even when the child's text is already in the parent's name.
This avoids hollowing out actual prose blocks where the \texttt{StaticText} child is the canonical carrier of the rendered text.

\paragraph{Icon characters.}
Before comparison, Unicode private-use characters in the range \texttt{U+E600--U+E6FF} (typically icon-font glyphs that render as visually empty placeholders) are stripped from both the parent name and the \texttt{StaticText} child.
Literal \texttt{\textbackslash uXXXX} escapes that appear in the serialized tree are also stripped.
This prevents icon characters from blocking an otherwise valid redundancy match.


\section{Effect of pruning on the augmented baselines}
\label{app:baseline-pruning}

\budgetactor{} applies rule-based accessibility-tree pruning, while AWM, ASI, and ReasoningBank are run in the released configurations of their original papers without pruning.
To characterize how the baselines behave under the same preprocessing, we apply \budgetactor{}'s pruning procedure to all three of them on the Shopping domain of WebArena, under Gemini~3~Flash and Qwen~3.6-27B.
Table~\ref{tab:baseline-pruning} reports the result.

\begin{table*}[ht]
\centering
\small
\begin{tabular}{lcccc}
\toprule
 & \multicolumn{2}{c}{\textbf{Gemini~3~Flash}} & \multicolumn{2}{c}{\textbf{Qwen~3.6-27B}} \\
\cmidrule(lr){2-3}\cmidrule(lr){4-5}
\textbf{Method} & \textbf{SR (\%) $\uparrow$} & \textbf{Tok (K) $\downarrow$} & \textbf{SR (\%) $\uparrow$} & \textbf{Tok (K) $\downarrow$} \\
\midrule
AWM              & 41.18{\scriptsize $\pm$2.1} & 70.8{\scriptsize $\pm$3.0} & 41.35{\scriptsize $\pm$1.9} & 71.8{\scriptsize $\pm$4.5} \\
\quad + pruning  & 43.32{\scriptsize $\pm$0.5} & 66.9{\scriptsize $\pm$4.3} & 39.04{\scriptsize $\pm$2.3} & 68.9{\scriptsize $\pm$2.7} \\
ASI              & 44.74{\scriptsize $\pm$1.2} & 82.8{\scriptsize $\pm$0.3} & 43.67{\scriptsize $\pm$3.5} & 78.8{\scriptsize $\pm$5.5} \\
\quad + pruning  & 45.63{\scriptsize $\pm$0.8} & 71.6{\scriptsize $\pm$5.6} & 45.81{\scriptsize $\pm$2.4} & 76.7{\scriptsize $\pm$6.3} \\
RBank            & 45.45{\scriptsize $\pm$1.4} & 54.6{\scriptsize $\pm$1.3} & 39.75{\scriptsize $\pm$0.8} & 65.4{\scriptsize $\pm$0.2} \\
\quad + pruning  & 42.25{\scriptsize $\pm$4.2} & 51.1{\scriptsize $\pm$0.2} & 39.21{\scriptsize $\pm$2.2} & 59.3{\scriptsize $\pm$1.0} \\
\midrule
\budgetactor{}   & \textbf{47.77}{\scriptsize $\pm$1.9} & \textbf{45.7}{\scriptsize $\pm$0.8} & \textbf{45.99}{\scriptsize $\pm$1.4} & \textbf{58.8}{\scriptsize $\pm$1.5} \\
\bottomrule
\end{tabular}
\caption{\textbf{Augmented baselines with and without \budgetactor{}'s pruning procedure} on the Shopping domain of WebArena, under Gemini~3~Flash and Qwen~3.6-27B. Each baseline is shown in its released configuration and with the same rule-based accessibility-tree pruning used by \budgetactor{} (\emph{+ pruning}). Values are means over three runs (mean~$\pm$~{\scriptsize std}). Success rate (\%) and total token usage per task (K).}
\label{tab:baseline-pruning}
\end{table*}

Pruning reduces token usage for every baseline under both models, consistent with its role as a cost-reduction mechanism.
Its effect on success rate is mixed: it helps AWM and ASI under Gemini~3~Flash and ASI under Qwen~3.6-27B, while slightly lowering the success rate of the remaining configurations.
Even where pruning helps, the pruned baseline still trails \budgetactor{} on success rate at a higher token cost.
The closest case is ASI under Qwen~3.6-27B, whose pruned variant reaches 45.81\% against 45.99\% for \budgetactor{}, but at 76.7K versus 58.8K tokens per task, roughly 30\% more.
Because pruning does not uniformly improve success rate and would depart from the configurations released by the original authors, we keep those released configurations as the primary baselines in the main results and report the pruned runs here.


\section{Results at a higher step budget}
\label{app:step-budget}

Our main experiments cap the augmented methods at a 10-step actor horizon and \budgetactor{} at 15 steps.
To check whether the accuracy--cost pattern persists at a higher operating point, we shift the maximum step count: the augmented methods from 10 to 15 and \budgetactor{} from 15 to 20.
We run this comparison under Qwen~3.6-27B on the Shopping, Reddit, and Admin domains of WebArena.
Table~\ref{tab:step-budget} reports the result.

\begin{table*}[ht]
\centering
\small
\begin{tabular}{lcccccc}
\toprule
 & \multicolumn{2}{c}{\textbf{Shopping}} & \multicolumn{2}{c}{\textbf{Reddit}} & \multicolumn{2}{c}{\textbf{Admin}} \\
\cmidrule(lr){2-3}\cmidrule(lr){4-5}\cmidrule(lr){6-7}
\textbf{Method} & \textbf{SR (\%) $\uparrow$} & \textbf{Tok (K) $\downarrow$} & \textbf{SR (\%) $\uparrow$} & \textbf{Tok (K) $\downarrow$} & \textbf{SR (\%) $\uparrow$} & \textbf{Tok (K) $\downarrow$} \\
\midrule
\budgetactor{} - 20 steps & \textbf{46.52}{\scriptsize $\pm$0.9} & \textbf{60.8}{\scriptsize $\pm$4.1} & \textbf{47.17}{\scriptsize $\pm$3.8} & \textbf{91.6}{\scriptsize $\pm$3.3} & 51.65{\scriptsize $\pm$2.9} & \textbf{161.4}{\scriptsize $\pm$22.4} \\
AWM - 15 steps   & 41.89{\scriptsize $\pm$3.4} & 79.5{\scriptsize $\pm$3.8} & 44.03{\scriptsize $\pm$0.5} & 105.0{\scriptsize $\pm$1.4} & 47.43{\scriptsize $\pm$4.2} & 212.6{\scriptsize $\pm$6.9} \\
ASI - 15 steps   & 45.63{\scriptsize $\pm$2.9} & 96.4{\scriptsize $\pm$1.4} & 43.08{\scriptsize $\pm$1.1} & 127.4{\scriptsize $\pm$13.7} & \textbf{51.83}{\scriptsize $\pm$3.0} & 206.8{\scriptsize $\pm$11.3} \\
RBank - 15 steps & 42.24{\scriptsize $\pm$2.8} & 79.1{\scriptsize $\pm$1.4} & 41.83{\scriptsize $\pm$2.7} & 111.6{\scriptsize $\pm$5.4} & 50.18{\scriptsize $\pm$1.1} & 198.0{\scriptsize $\pm$24.2} \\
\bottomrule
\end{tabular}
\caption{\textbf{Success rate (\%) and total token usage per task (K) at a higher step budget}, under Qwen~3.6-27B on WebArena. The augmented methods use a 15-step actor horizon and \budgetactor{} a 20-step horizon. Values are means over three runs (mean~$\pm$~{\scriptsize std}).}
\label{tab:step-budget}
\end{table*}

The ordering is largely preserved at this second operating point.
\budgetactor{} has the highest success rate on Shopping and Reddit and uses fewer tokens than every augmented method in all three domains.
On Admin, ASI is nominally $0.18$ points above \budgetactor{} (51.83\% vs.\ 51.65\%), a gap far smaller than the run-level standard deviations, and it spends 28\% more tokens to reach it (206.8K vs.\ 161.4K).
Averaged over the three domains, \budgetactor{} leads ASI on both axes, at 48.6\% success and 106.2K tokens per task versus 47.4\% and 145.6K.
Granting the augmented methods more steps therefore does not let them overtake the vanilla actor: they improve, but so does \budgetactor{} under its matched increase, and both the ordering and the cost gap remain intact.


\section{Prompt templates}
\label{app:prompts}

\budgetactor{}, AWM, and ASI use the prompt templates released by \citet{wang-etal-2025-asi} without modification, except as noted below.
Since the templates are publicly available in the ASI code repository, we do not report them here.

\paragraph{Model-specific modification.}
The prompts are model-agnostic with one exception.
When running experiments with GPT-5.4-mini, we observed two behaviors that interfered with evaluation.
First, the model frequently produced action outputs with no accompanying reasoning or explanation.
Such explanations are important for augmentation methods: the induction modules parse them to extract reusable knowledge from trajectories.
Second, the model sometimes generated turns asking the user for clarification or additional input, which is not possible during automated benchmark evaluation.
% 
To address both issues, the following text was appended to the action prompt for GPT-5.4-mini:

\begin{quote}
\itshape
You cannot ask for follow-up questions or hand back to the user.
Do your best with the information already available.
Output a summary of your reasoning and the important information you have found, followed by the best actions at this step (in one triple backticks block).
\end{quote}

\noindent It is worth noting that \budgetactor{} is less dependent on reasoning traces in the output: the vanilla agent acts only on the current observation, so in many cases explanatory text is overhead that could in principle be removed to reduce completion tokens.
This represents an additional efficiency advantage of \budgetactor{} that our numbers do not capture.
However, to ensure a fair comparison across all methods, we applied the same prompt to all methods under each model, including the above addition for GPT-5.4-mini.


\section{Failure mode analysis of augmented agents}
\label{app:failure-modes}

Section~\ref{sec:fragility} discusses the structural fragility of programmatic skills in dynamic accessibility-tree interfaces at a conceptual level.
This appendix extends that discussion with quantitative characterizations of failure modes across all three augmented methods, and presents concrete examples of tasks where augmented agents fail while \budgetactor{} succeeds.
The analysis in this section covers the Shopping, Reddit, and Admin domains under Gemini~3~Flash and GPT-5.4-mini.
Unless stated otherwise, quantities are averaged over three independent runs per domain--model pair.

\subsection{ASI: Verification confounds function success with agent recovery}
\label{app:asi-failures}

ASI induces Python functions from successful trajectories and adds them to a shared action library after a verification episode.
The intended guarantee is that only working functions enter the library.
In the verification procedure, the induced function is called as the first decision. If the function changes the environment but fails on one of its underlying actions, the agent may still recover via primitive actions and solve the task. In that case, the episode can be judged a success and the function is written to disk, without actual verification.
A common instance of this failure is when the function fails because an element ID has changed after the page reloaded for the verification episode.
Functions that require arguments whose values depend on a prior action's page state are especially prone to this: the caller cannot supply a valid BID for an element that only becomes visible after an earlier click.
The induced functions are also typically short, averaging approximately 2.4 primitive actions per function under Gemini~3~Flash and 2.5 under GPT-5.4-mini, which limits the abstraction benefit.
Some functions also hardcode task-specific values that do not generalize across tasks; for example, the function \texttt{fill\_refund\_request\_form} in the shopping library hardcodes the complaint message ``\textit{It broke after just three days of use}'' regardless of the actual reason for the refund.

Table~\ref{tab:asi-verification} quantifies the verification-confound effect.
We define a \emph{first-step verification failure} as a verification run in which the induced function fails before the second agent action.
A \emph{recovered} verification is one that was still judged correct despite that first-step failure.
First-step failures are not rare.
Under GPT, the induced function fails immediately in 72.2\% of Shopping verifications; over half of those are recovered and still produce a stored library entry.
Even under Gemini Shopping, where the first-step failure rate is lower (33.3\%), two thirds of failures recover.
In both cases the result is a function library that contains entries whose correctness at verification time was not due to the function itself.
\\

\begin{table}[ht]
\centering
\resizebox{\linewidth}{!}{
\begin{tabular}{lccc}
\toprule
\textbf{Domain} & \textbf{Attempts} & \textbf{First-step fail} & \textbf{Recov.~rate} \\
\midrule
\multicolumn{4}{c}{\textbf{Gemini~3~Flash}} \\
\midrule
Shopping & 58.0 & 19.3 & 69.0\% \\
Reddit   & 30.0 &  9.0 & 59.3\% \\
Admin    & 64.3 &  6.3 & 63.2\% \\
\midrule
\multicolumn{4}{c}{\textbf{GPT-5.4-mini}} \\
\midrule
Shopping & 44.3 & 32.0 & 54.2\% \\
Reddit   & 25.0 & 12.3 & 45.9\% \\
Admin    & 51.7 & 22.7 & 47.1\% \\
\bottomrule
\end{tabular}
}
\caption{\textbf{ASI verification failures at the first high-level function call} on WebArena. Values are averages over three runs. \emph{Attempts}: total number of function induction attempts by ASI across all tasks. \emph{First-step fail}: number of attempts where the induced function failed before the second agent action. \emph{Recov.~rate}: fraction of first-step failures where the verification episode was nonetheless judged correct, causing a potentially broken function to be stored in the shared library.}
\label{tab:asi-verification}
\end{table}

\subsection{AWM: Generic templates and induction from failed trajectories}
\label{app:awm-failures}

AWM induces abstract natural-language workflows from trajectories the judge deems successful, then retrieves and injects the closest workflow into the actor's prompt at each inference step.
Many induced workflows are too generic to add navigational value.
A stored shopping workflow for order-history aggregation reduces to: ``\textit{click the account link, click the full-order-history link, send a summary}''.
This describes the high-level intent without preserving operational details that determine correctness, such as whether pagination is needed, what the stopping condition is, or how totals are computed.
Similarly, a stored price-range workflow directs the agent to sort, toggle direction, navigate to the last page, and report the range, but omits checks for sort direction availability, filtered result sets, or whether the last-page element is reachable in the current state.
Because the workflow is injected into the actor's prompt, the agent must decide at each step whether to follow it or ignore it, a burden that smaller models are especially ill-equipped to handle, as it competes with the current accessibility-tree and task instruction.

A more fundamental problem is that AWM's induction is gated only by the judge's verdict.
Table~\ref{tab:awm-induction} shows the fraction of workflow-induction events triggered by failed trajectories.
On Shopping, nearly half of all induction events, 49.5\% under Gemini and 52.3\% under GPT, originate from failed agent trajectories.
The contamination is substantially lower for Reddit (11--19\%), which may help explain AWM's stronger performance on that domain.
Because AWM's duplicate-suppression mechanism checks only whether a new workflow semantically overlaps an existing one, it does not filter out strategies that are simply wrong: once an incorrect workflow enters the library, it can be retrieved and applied to all future tasks in the same category.

\begin{table}[ht]
\centering
\resizebox{\linewidth}{!}{
\begin{tabular}{lccc}
\toprule
\textbf{Domain} & \textbf{Induced} & \textbf{From failed (\%)} & \textbf{Final WFs} \\
\midrule
\multicolumn{4}{c}{\textbf{Gemini~3~Flash}} \\
\midrule
Shopping & 105.7 & 49.5\% & 37.3 \\
Reddit   &  45.7 & 10.9\% & 16.7 \\
Admin    &  99.7 & 42.1\% & 48.3 \\
\midrule
\multicolumn{4}{c}{\textbf{GPT-5.4-mini}} \\
\midrule
Shopping &  78.3 & 52.3\% & 33.7 \\
Reddit   &  31.3 & 19.1\% & 13.7 \\
Admin    &  50.0 & 42.0\% & 24.3 \\
\bottomrule
\end{tabular}
}
\caption{\textbf{AWM workflow-induction statistics} on WebArena, averaged per run. \emph{Induced}: total induction events triggered by the judge for a run. \emph{From failed (\%)}: fraction originating from tasks the ground-truth evaluator classified as failed. \emph{Final WFs}: distinct workflows in the final library after duplicate suppression.}
\label{tab:awm-induction}
\end{table}

\subsection{ReasoningBank: Contamination and memory over-application}
\label{app:rb-failures}

ReasoningBank stores a few short lessons from every task trajectory regardless of outcome, labeling each entry as a success or failure based on the judge's verdict.
The $K$ most semantically similar entries are retrieved at inference time and injected into the actor's context.
Table~\ref{tab:rb-corpus} reports the size and contamination of the resulting memory corpus.
On Shopping tasks, more than half of all success-labeled entries, 52.9\% for Gemini and 59.5\% for GPT, come from trajectories that failed according to the ground-truth evaluator.
These entries are framed with positive guidance language regardless of their true provenance, so the retrieval mechanism cannot distinguish between lessons that worked and lessons that did not.

\begin{table}[ht]
\centering
\resizebox{\linewidth}{!}{%
\begin{tabular}{lcccc}
\toprule
\textbf{Domain} & \shortstack{\textbf{Total}\\\textbf{mem.}} & \shortstack{\textbf{Succ.-}\\\textbf{labeled}} & \shortstack{\textbf{FP in}\\\textbf{succ.~(\%)}} & \shortstack{\textbf{FN in}\\\textbf{fail~(\%)}} \\
\midrule
\multicolumn{5}{c}{\textbf{Gemini~3~Flash}} \\
\midrule
Shopping & 518 & 96 & 52.9\% & 40.7\% \\
Reddit   & 294 & 32 & 30.2\% & 25.3\% \\
Admin    & 507 & 86 & 42.6\% & 40.4\% \\
\midrule
\multicolumn{5}{c}{\textbf{GPT-5.4-mini}} \\
\midrule
Shopping & 544 & 67 & 59.5\% & 18.4\% \\
Reddit   & 307 & 20 & 60.0\% & 17.4\% \\
Admin    & 532 & 62 & 48.1\% & 25.2\% \\
\bottomrule
\end{tabular}}
\caption{\textbf{ReasoningBank memory-corpus statistics} on WebArena, averaged over three runs. Tasks processed per run: 187 (Shopping), 106 (Reddit), 182 (Admin). \emph{Succ.-labeled}: tasks labeled successful by the judge. \emph{FP in succ.}: fraction of success-labeled tasks from trajectories that failed according to the ground-truth evaluator. \emph{FN in fail}: fraction of fail-labeled tasks from trajectories that succeeded according to the ground-truth evaluator.}

\label{tab:rb-corpus}
\end{table}

A separate concern is injection pressure.
Even though only a few memories are retrieved per task, they are referenced heavily in the agent traces: across the domain--model configurations in Table~\ref{tab:rb-corpus}, we observe an average of 8--12 explicit memory-item mentions per task, with over 95\% of Gemini tasks and over 75\% of GPT tasks referencing at least one retrieved memory.
The retrieved memories therefore function as active policy hints rather than passive background.
When a memory is overly general or mildly mismatched to the current task, the agent may prioritize it over the current accessibility-tree, producing extra navigation steps, unnecessary checks, or incorrect early termination.
\budgetactor{}, which carries no retrieved memory layer, avoids this class of distraction entirely by grounding each action in the current observation.

\subsection{Discussion: Structural limits of judge-based quality control}
\label{app:judge-limits}

The contamination patterns described above for AWM and ReasoningBank, and the verification confound in ASI, share a common structural root: each method delegates quality control to an LLM judge whose input is fundamentally incomplete for the tasks it is asked to evaluate.

The judge receives the agent's textual output (its reasoning trace and final response) and, in some configurations, a snapshot of the browser state at the end of the episode.
It does not observe the sequence of intermediate page states the agent traversed, the specific elements it interacted with, or whether information the agent reports was actually encountered during navigation or confabulated in the reasoning trace.
For tasks where correctness depends on having visited the right pages and extracted the right values, a judge working from this terminal view cannot reliably reconstruct whether the work was actually done.

This is not necessarily a matter of judge capability.
A more powerful model receiving the same partial input faces the same gap: the evidence needed to verify navigation-dependent correctness is simply absent from the judge's context.
Nor is it straightforwardly fixable by expanding what is passed to the judge.
Supplying the full trace at every step
would require passing on the order of hundreds of thousands of tokens per judge call for a ten-step task on a content-heavy page, making each evaluation call more expensive than the agent trajectory it is judging, and possibly exceeding practical context limits entirely.

\begin{table*}[t!]
\centering
\small
\resizebox{\textwidth}{!}{%
\begin{tabular}{lp{3.0cm}p{5.6cm}p{5.6cm}}
\toprule
\textbf{Method} & \textbf{Intent} & \textbf{Augmented agent behavior} & \textbf{\budgetactor{} behavior} \\
\midrule
\addlinespace
AWM
  & ``Buy the highest-rated product in Ceiling Light within a budget above \$1000.''
  & Retrieved workflow for ``Buy the highest-rated product'' prescribes sorting then clicking the top result. Agent buys the wrong product because the workflow does not account for rating ties or the budget threshold.
  & Applies sorting, checks rating and price against the budget, and selects the correct product. \\
\addlinespace
ASI
  & ``Change the delivery address for my most recent order to 3 Oxford St, Cambridge, MA.''
  & Calls \texttt{navigate\_to\_contact\_page} and \texttt{fill\_contact\_form} using the element ID from the docstring example (\texttt{'1381'}). Receives \texttt{ValueError} and \texttt{TypeError}; form submission fails.
  & Navigates to Address Book, adds the new address, recognizes that the order address cannot be changed through the UI, and reports accordingly. \\
\addlinespace
RBank
  & ``What are the top-5 best-selling products in 2023?''
  & Retrieved memory on handling ranking ties directs attention to tie-breaking. Agent misidentifies the 4th and 5th items (reports \emph{Sparta Gym Tank} and \emph{Angel Light Running Short} instead of \emph{Sprite Stasis Ball} and \emph{Hawkeye Yoga Short}).
  & Reads the same report without distraction and returns the correct ranked list. \\
\bottomrule
\end{tabular}}
\caption{\textbf{Example failures of augmented agents} where \budgetactor{} succeeds. Tasks are selected from WebArena (single runs shown; all from Gemini~3~Flash).}
\label{tab:failure-examples}
\end{table*}

\subsection{Illustrative task-level failures}
\label{app:failure-examples}

Table~\ref{tab:failure-examples} presents representative task-level failures drawn from our Gemini~3~Flash runs, showing in each case what the augmented agent does and how \budgetactor{} succeeds on the same task.


\section{Per-run variance results}
\label{app:variance}

Tables~\ref{tab:variance-gemini}, \ref{tab:variance-gpt}, and~\ref{tab:variance-qwen} report success rate and token usage for each method across three independent runs under Gemini~3~Flash, GPT-5.4-mini, and Qwen~3.6-27B, respectively.
We report sample standard deviations across the three runs.

\begin{table*}[ht]
\centering
\small
\begin{tabular}{lcccccccc}
\toprule
& \multicolumn{4}{c}{\textbf{Success Rate (\%) $\uparrow$}} & \multicolumn{4}{c}{\textbf{Tok/task (K) $\downarrow$}} \\
\cmidrule(lr){2-5} \cmidrule(lr){6-9}
\textbf{Method} & \textbf{Run a} & \textbf{Run b} & \textbf{Run c} & \textbf{Mean$\pm$Std} & \textbf{Run a} & \textbf{Run b} & \textbf{Run c} & \textbf{Mean$\pm$Std} \\
\midrule
\multicolumn{9}{c}{\textbf{Shopping}} \\
\midrule
\budgetactor{}      & 47.59 & 49.73 & 45.99 & \textbf{47.77}{\scriptsize $\pm$1.9} & 46.6 & 45.5 & 45.1 & \textbf{45.7}{\scriptsize $\pm$0.8} \\
AWM            & 43.32 & 39.04 & 41.18 & 41.18{\scriptsize $\pm$2.1} & 68.0 & 70.6 & 73.9 & 70.8{\scriptsize $\pm$3.0} \\
ASI            & 43.32 & 45.45 & 45.45 & 44.74{\scriptsize $\pm$1.2} & 82.4 & 83.0 & 83.0 & 82.8{\scriptsize $\pm$0.3} \\
RBank          & 44.39 & 44.91 & 47.05 & 45.45{\scriptsize $\pm$1.4} & 56.0 & 54.1 & 53.6 & 54.6{\scriptsize $\pm$1.3} \\
\midrule
\multicolumn{9}{c}{\textbf{Reddit}} \\
\midrule
\budgetactor{}      & 48.11 & 47.17 & 47.17 & \textbf{47.48}{\scriptsize $\pm$0.5} & 71.7 & 69.5 & 73.8 & \textbf{71.7}{\scriptsize $\pm$2.2} \\
AWM            & 49.06 & 46.23 & 47.17 & \textbf{47.48}{\scriptsize $\pm$1.4} & 87.6 & 88.9 & 85.0 & 87.2{\scriptsize $\pm$2.0} \\
ASI            & 47.17 & 43.40 & 44.34 & 44.97{\scriptsize $\pm$2.0} & 98.5 & 89.3 & 96.0 & 94.6{\scriptsize $\pm$4.8} \\
RBank          & 39.62 & 38.68 & 41.50 & 39.93{\scriptsize $\pm$1.4} & 78.9 & 73.6 & 77.8 & 76.8{\scriptsize $\pm$2.8} \\
\midrule
\multicolumn{9}{c}{\textbf{Admin}} \\
\midrule
\budgetactor{}      & 53.85 & 56.59 & 56.59 & \textbf{55.68}{\scriptsize $\pm$1.6} &  95.2 &  98.6 & 103.2 &  \textbf{99.0}{\scriptsize $\pm$4.0} \\
AWM            & 47.80 & 47.80 & 46.70 & 47.43{\scriptsize $\pm$0.6} & 136.6 & 147.0 & 144.4 & 142.7{\scriptsize $\pm$5.4} \\
ASI            & 53.85 & 51.65 & 52.75 & 52.75{\scriptsize $\pm$1.1} & 132.0 & 152.3 & 133.9 & 139.4{\scriptsize $\pm$11.2} \\
RBank          & 47.80 & 51.10 & 47.80 & 48.90{\scriptsize $\pm$1.9} & 120.8 & 126.4 & 126.5 & 124.6{\scriptsize $\pm$3.3} \\
\midrule
\multicolumn{9}{c}{\textbf{GitLab}} \\
\midrule
\budgetactor{}      & 28.33 & 29.44 & 29.44 & \textbf{29.07}{\scriptsize $\pm$0.6} & 79.8 & 76.2 & 78.0 & 78.0{\scriptsize $\pm$1.8} \\
AWM            & 25.00 & 23.33 & 25.00 & 24.44{\scriptsize $\pm$1.0} & 93.1 & 91.0 & 92.6 & 92.3{\scriptsize $\pm$1.1} \\
ASI            & 22.78 & 21.11 & 25.00 & 22.96{\scriptsize $\pm$2.0} & 102.3 & 107.0 & 113.8 & 107.7{\scriptsize $\pm$5.8} \\
RBank          & 22.78 & 23.89 & 22.22 & 22.96{\scriptsize $\pm$0.9} & 71.9 & 72.4 & 73.8 & \textbf{72.7}{\scriptsize $\pm$1.0} \\
\bottomrule
\end{tabular}
\caption{\textbf{Per-run success rate (\%) and token usage per task (K)} across three independent runs under Gemini~3~Flash on WebArena.}
\label{tab:variance-gemini}
\end{table*}

\begin{table*}[ht]
\centering
\small
\begin{tabular}{lcccccccc}
\toprule
& \multicolumn{4}{c}{\textbf{Success Rate (\%) $\uparrow$}} & \multicolumn{4}{c}{\textbf{Tok/task (K) $\downarrow$}} \\
\cmidrule(lr){2-5} \cmidrule(lr){6-9}
\textbf{Method} & \textbf{Run a} & \textbf{Run b} & \textbf{Run c} & \textbf{Mean$\pm$Std} & \textbf{Run a} & \textbf{Run b} & \textbf{Run c} & \textbf{Mean$\pm$Std} \\
\midrule
\multicolumn{9}{c}{\textbf{Shopping}} \\
\midrule
\budgetactor{}      & 37.43 & 39.04 & 39.04 & \textbf{38.50}{\scriptsize $\pm$0.9} & 55.1 & 53.5 & 54.2 & \textbf{54.3}{\scriptsize $\pm$0.8} \\
AWM            & 27.81 & 31.02 & 31.02 & 29.95{\scriptsize $\pm$1.9} & 64.2 & 61.6 & 60.2 & 62.0{\scriptsize $\pm$2.0} \\
ASI            & 34.22 & 32.62 & 32.62 & 33.15{\scriptsize $\pm$0.9} & 70.6 & 83.5 & 66.2 & 73.4{\scriptsize $\pm$9.0} \\
RBank          & 21.93 & 27.81 & 29.41 & 26.38{\scriptsize $\pm$3.9} & 59.3 & 58.2 & 56.2 & 57.9{\scriptsize $\pm$1.6} \\
\midrule
\multicolumn{9}{c}{\textbf{Reddit}} \\
\midrule
\budgetactor{}      & 38.68 & 35.85 & 29.25 & \textbf{34.59}{\scriptsize $\pm$4.8} & 89.4 & 102.5 & 96.6 & 96.2{\scriptsize $\pm$6.6} \\
AWM            & 30.19 & 30.19 & 28.30 & 29.56{\scriptsize $\pm$1.1} & 77.1 & 94.4 & 69.8 & 80.4{\scriptsize $\pm$12.6} \\
ASI            & 26.42 & 30.19 & 30.19 & 28.93{\scriptsize $\pm$2.2} & 83.8 & 85.0 & 92.8 & 87.2{\scriptsize $\pm$4.9} \\
RBank          & 18.87 & 24.53 & 22.64 & 22.01{\scriptsize $\pm$2.9} & 67.8 & 72.5 & 67.3 & \textbf{69.2}{\scriptsize $\pm$2.9} \\
\midrule
\multicolumn{9}{c}{\textbf{Admin}} \\
\midrule
\budgetactor{}      & 35.16 & 35.16 & 37.36 & \textbf{35.89}{\scriptsize $\pm$1.3} & 117.4 & 133.0 & 121.9 & 124.1{\scriptsize $\pm$8.0} \\
AWM            & 28.57 & 33.52 & 34.62 & 32.24{\scriptsize $\pm$3.2} & 129.3 &  95.9 & 114.9 & \textbf{113.4}{\scriptsize $\pm$16.8} \\
ASI            & 28.57 & 35.16 & 35.16 & 32.96{\scriptsize $\pm$3.8} & 126.1 & 146.5 & 142.9 & 138.5{\scriptsize $\pm$10.9} \\
RBank          & 36.26 & 32.97 & 33.52 & 34.25{\scriptsize $\pm$1.8} & 121.0 & 125.5 & 122.8 & 123.1{\scriptsize $\pm$2.3} \\
\midrule
\multicolumn{9}{c}{\textbf{GitLab}} \\
\midrule
\budgetactor{}      & 20.00 & 24.44 & 22.22 & \textbf{22.22}{\scriptsize $\pm$2.2} & 94.4 & 85.5 & 89.6 & 89.8{\scriptsize $\pm$4.4} \\
AWM            & 16.11 & 18.33 & 17.22 & 17.22{\scriptsize $\pm$1.1} & 92.5 & 97.0 & 97.5 & 95.7{\scriptsize $\pm$2.7} \\
ASI            & 21.11 & 21.11 & 20.00 & 20.74{\scriptsize $\pm$0.6} & 94.8 & 97.9 & 94.4 & 95.7{\scriptsize $\pm$1.9} \\
RBank          & 11.11 & 15.56 & 16.67 & 14.45{\scriptsize $\pm$2.9} & 70.9 & 70.3 & 66.8 & \textbf{69.3}{\scriptsize $\pm$2.2} \\
\bottomrule
\end{tabular}
\caption{\textbf{Per-run success rate (\%) and token usage per task (K)} across three independent runs under GPT-5.4-mini on WebArena.}
\label{tab:variance-gpt}
\end{table*}

\begin{table*}[ht]
\centering
\small
\begin{tabular}{lcccccccc}
\toprule
& \multicolumn{4}{c}{\textbf{Success Rate (\%) $\uparrow$}} & \multicolumn{4}{c}{\textbf{Tok/task (K) $\downarrow$}} \\
\cmidrule(lr){2-5} \cmidrule(lr){6-9}
\textbf{Method} & \textbf{Run a} & \textbf{Run b} & \textbf{Run c} & \textbf{Mean$\pm$Std} & \textbf{Run a} & \textbf{Run b} & \textbf{Run c} & \textbf{Mean$\pm$Std} \\
\midrule
\multicolumn{9}{c}{\textbf{Shopping}} \\
\midrule
\budgetactor{}      & 47.06 & 44.39 & 46.52 & \textbf{45.99}{\scriptsize $\pm$1.4} & 57.9 & 60.5 & 57.9 & \textbf{58.8}{\scriptsize $\pm$1.5} \\
AWM            & 39.57 & 43.30 & 41.18 & 41.35{\scriptsize $\pm$1.9} & 73.6 & 66.7 & 75.1 & 71.8{\scriptsize $\pm$4.5} \\
ASI            & 47.06 & 43.85 & 40.10 & 43.67{\scriptsize $\pm$3.5} & 72.6 & 81.1 & 82.8 & 78.8{\scriptsize $\pm$5.5} \\
RBank          & 39.57 & 40.64 & 39.04 & 39.75{\scriptsize $\pm$0.8} & 65.6 & 65.2 & 65.3 & 65.4{\scriptsize $\pm$0.2} \\
\midrule
\multicolumn{9}{c}{\textbf{Reddit}} \\
\midrule
\budgetactor{}      & 47.17 & 43.40 & 42.45 & \textbf{44.34}{\scriptsize $\pm$2.5} & 94.5 & 86.5 & 88.0 & \textbf{89.7}{\scriptsize $\pm$4.3} \\
AWM            & 44.34 & 45.28 & 39.62 & 43.08{\scriptsize $\pm$3.0} & 110.3 &  87.8 & 112.8 & 103.6{\scriptsize $\pm$13.8} \\
ASI            & 43.40 & 41.51 & 44.34 & 43.08{\scriptsize $\pm$1.4} & 118.1 & 125.9 & 117.6 & 120.5{\scriptsize $\pm$4.7} \\
RBank          & 39.62 & 41.51 & 42.45 & 41.19{\scriptsize $\pm$1.4} &  96.1 &  95.9 & 105.4 &  99.1{\scriptsize $\pm$5.4} \\
\midrule
\multicolumn{9}{c}{\textbf{Admin}} \\
\midrule
\budgetactor{}      & 49.45 & 52.20 & 50.54 & \textbf{50.73}{\scriptsize $\pm$1.4} & 141.2 & 125.3 & 124.4 & \textbf{130.3}{\scriptsize $\pm$9.5} \\
AWM            & 45.60 & 43.40 & 49.45 & 46.15{\scriptsize $\pm$3.1} & 192.8 & 172.3 & 170.3 & 178.5{\scriptsize $\pm$12.5} \\
ASI            & 51.65 & 48.35 & 47.25 & 49.08{\scriptsize $\pm$2.3} & 177.7 & 183.1 & 185.9 & 182.2{\scriptsize $\pm$4.2} \\
RBank          & 46.15 & 47.25 & 49.45 & 47.62{\scriptsize $\pm$1.7} & 160.8 & 162.9 & 147.1 & 156.9{\scriptsize $\pm$8.6} \\
\midrule
\multicolumn{9}{c}{\textbf{GitLab}} \\
\midrule
\budgetactor{}      & 27.22 & 29.44 & 27.78 & \textbf{28.15}{\scriptsize $\pm$1.2} & 109.0 & 96.2 & 100.2 & 101.8{\scriptsize $\pm$6.5} \\
AWM            & 23.33 & 22.22 & 24.44 & 23.33{\scriptsize $\pm$1.1} & 102.9 & 106.6 & 97.0 & 102.2{\scriptsize $\pm$4.8} \\
ASI            & 27.22 & 25.55 & 24.44 & 25.74{\scriptsize $\pm$1.4} & 120.4 & 117.0 & 124.3 & 120.6{\scriptsize $\pm$3.6} \\
RBank          & 20.56 & 22.78 & 19.44 & 20.93{\scriptsize $\pm$1.7} & 81.2 & 84.6 & 92.3 & \textbf{86.0}{\scriptsize $\pm$5.7} \\
\bottomrule
\end{tabular}
\caption{\textbf{Per-run success rate (\%) and token usage per task (K)} across three independent runs under Qwen~3.6-27B on WebArena.}
\label{tab:variance-qwen}
\end{table*}
